% ============================================================================
\section{Evolución del núcleo}
\label{sec:evolucion}
% ============================================================================

Cada sección cerró con un recuadro de evolución propio. Esta sección los
ordena en el tiempo y extrae el patrón común. El sistema vigente reemplazó a
un conjunto de repositorios anteriores. No fue una reescritura: lo probado se
portó y lo que cambió fue el modelo de dominio, que era donde estaban los
defectos. Todas las fechas son de 2026.

\begin{figure}[htbp]
\centering
\begin{tikzpicture}[
  hito/.style={circle, fill=azul, inner sep=1.6pt},
  txt/.style={font=\footnotesize, align=left, text width=11.2cm, anchor=west},
  fecha/.style={font=\footnotesize\bfseries, anchor=east},
]
\draw[thick, gris] (0,0.5) -- (0,-9.6);
\foreach \y/\f/\t in {
  0/{21 sep}/{Monorrepositorio nuevo. Se portan el catálogo y el ciclo editorial; el modelo de dominio se rehace (cantidad y empaque separados, estados \textsc{SinMatch} y \textsc{HuecoCatálogo} distintos).},
  -1.5/{22 sep}/{Ciclo completo de punta a punta. Percepción todo-o-nada con rechazo por grados; normalización; emparejador léxico con cabeza anclada; optimizador \textbf{v1} (heurística declarada como andamio).},
  -3.4/{23 sep}/{Resolución híbrida: canal semántico con centroides de concepto y fusión por guardas. Optimizador \textbf{v2} (NSGA-II, variables reales y piso) y, el mismo día, \textbf{v3} (codificación entera, reparación por redondeo, población y generaciones escaladas, la biblioteca de optimización como dependencia obligatoria).},
  -5.6/{4 oct}/{\textbf{Evaluación} contra el óptimo exacto y un modelo de propósito general. Optimizador \textbf{v4} (costo de línea, siembra con la escalera, mejor marca que cabe) y \textbf{v5} (kits a pedido, A~tu~Medida, tope sin estirar). Curaduría de las 13 listas de un colegio: correcciones al emparejador caso por caso, concepto asignado que manda, tipo del concepto, reglas de segmentación y reintento en la percepción.},
  -8.6/{7 oct}/{Optimizador \textbf{v6}: herramientas por hijo. Vectores recalculados con huella nueva. Atributos del catálogo que filtran por lo que el colegio pide. Tras la formalización de este documento, optimizador \textbf{v7} (proyección sobre lo surtible) y coseno recortado en 0 en la resolución.}
}{
  \node[hito] at (0,\y) {};
  \node[fecha] at (-0.25,\y) {\f};
  \node[txt] at (0.3,\y) {\t};
}
\end{tikzpicture}
\caption{Línea de tiempo del núcleo algorítmico.}
\label{fig:linea-tiempo}
\end{figure}

\subsection{El patrón: defecto medido, cambio mínimo, efecto medido}

Casi todos los cambios del núcleo siguen el mismo ciclo:
\begin{enumerate}
  \item \textbf{Un defecto observable en un caso real}, no hipotético: un
        renglón concreto que se clasificaba mal, una lista concreta que se
        rechazaba, una cifra concreta de la evaluación.
  \item \textbf{Un diagnóstico de la causa}, distinto del síntoma. Por
        ejemplo, «lápices de colores» caía en esfero; el síntoma sugería
        alargar el prefijo comparado, pero la causa era comparar formas
        flexionadas, y la corrección fue singularizar.
  \item \textbf{El cambio más pequeño que elimina la causa}, fijado por una
        prueba con el caso original.
  \item \textbf{La verificación de que los vecinos no se mueven}: al cambiar
        un sinónimo o una regla, se vuelve a emparejar el corpus para comprobar
        que ningún otro renglón cambió de concepto.
\end{enumerate}

\Cref{tab:evolucion-resumen} resume los cambios de mayor efecto.

\begin{table}[htbp]
\centering\small
\caption{Cambios de mayor efecto en el núcleo.}
\label{tab:evolucion-resumen}
\begin{tabular}{@{}L{2.2cm}L{4.6cm}L{4.2cm}L{3.6cm}@{}}
\toprule
Etapa & Antes & Después & Evidencia \\
\midrule
Percepción & tope de 5 páginas que recortaba en silencio & todo o nada; guarda de 40 páginas; rechazo por grados & kits incompletos presentados como completos \\
Percepción & página que continúa marcada como incompleta & regla de continuación & 2 listas reales rechazadas \\
Percepción & un fallo de red invalidaba la lista & un reintento de transporte & página colgada 60~s, 6~s al reintentar \\
Normalización & cantidad y empaque en un número & $q$ y $\emp$ separados, $\techo{q/\emp}$ & «doce cajas» por doce lápices \\
Resolución & contención en cualquier posición & cabeza anclada & chaquiras $\to$ lápices de colores \\
Resolución & solo léxico & fusión con canal semántico & sinónimos no escritos («estuche») \\
Resolución & canal semántico inactivo en producción & una sola carga del índice & coseno siempre 0 \\
Catálogo & empaque 1 en todos los productos & regla «x$N$» salvo medida & la fórmula $\techo{q/\emp}$ no hacía nada \\
Catálogo & escaleras de nivel invertidas & nivelación por permutación & Ahorrador más caro que Premium \\
Optimización & NSGA-II separable que empataba con fuerza bruta & enumeración exacta + NSGA-II sembrado & pruebas del prototipo \\
Optimización & índice real con piso & índice entero con redondeo & SBX interpolaba etiquetas \\
Optimización & orden por precio unitario & orden por costo de línea & Ahorrador 80\,\% $\to$ 100\,\% \\
Optimización & recomendado con presupuesto = más barato & mejor marca que cabe, del frente & calidad 40\,\% $\to$ óptimo \\
Optimización & tres kits siempre & kits a pedido; A~tu~Medida & kits que nadie pedía \\
Optimización & herramientas cobradas por materia & una por hijo (máximo, no suma) & 43 renglones repetidos \\
Optimización & un candidato sin precio parecía gratis en el frente & proyección sobre los puntos de máxima disponibilidad & instancia de dos líneas (defecto latente) \\
Resolución & un coseno negativo podía restar & $\kappa$ recortado en 0 & propiedad «híbrido $\geq$ léxico» \\
\bottomrule
\end{tabular}
\end{table}

\subsection{Lo que se decidió no cambiar}

La evolución también se ve en lo que se conservó a propósito. La aritmética
de unidades, la regla de que lo informativo no se cobra y la separación entre
\textsc{SinMatch} y \textsc{HuecoCatálogo} se definieron en la primera
versión y no cambiaron en ninguna de las siguientes, porque no dependen del
algoritmo. El optimizador se reemplazó tres veces sin tocarlas. Esa
estabilidad es el resultado del principio «separar decidir de presentar»
(\cref{sec:principios}).
